\section{An induction over the cyclic cut order}
\label{sec:path-fixed-cut-induction}

Deleting one edge of a cyclic cut reduces its order by one after two suppressions.
The uniform terminal-path bound controls the cycles lost in this deletion.
This gives the sharp error for every fixed cyclic cut order.

Write $\lambda_c(G)$ for the minimum size of an edge cut whose two shores both contain cycles.
Set $\lambda_c(G)=\infty$ when no such cut exists.

\begin{lemma}[Smoothing one edge of a minimum cyclic cut]
\label{lem:path-minimum-cyclic-cut-smoothing}
Let $G$ be a connected cubic multigraph of girth at least five, and suppose $\lambda_c(G)=k\ge5$.
Choose a minimum cyclic cut $\delta_G A$, an edge $e=uv$ in it, and put $N=G-e$.
Suppress the degree-two vertices $u,v$ in $N$, obtaining $G'$.
Then
\begin{enumerate}
\item $G'$ is connected and loopless cubic, with $\beta(G')=\beta(G)-1$;
\item the two original shores give a cyclic $(k-1)$-cut in $G'$, with unchanged shore ranks;
\item $\lambda_c(G')=k-1$;
\item $c(G)=c(G')+P_{uv}(N)$.
\end{enumerate}
\end{lemma}

\begin{proof}
Each shore of the minimum cut is connected.
If one shore were disconnected, a cyclic component on that shore would have boundary smaller than $k$.
Its complement would still contain the opposite shore's cycle.

Each vertex has at most one incidence in the cut.
Otherwise its internal degree is at most one, so it lies on no cycle in its shore.
Moving it across preserves a cycle on each side and decreases the cut size.
Thus $u$ has two neighbors inside $A$, and $v$ has two inside its opposite shore.
The girth assumption makes each pair of internal neighbors distinct.

Every bridge in a cubic multigraph is a cyclic cut.
Indeed, a bridge shore with $n$ vertices and rank $h$ satisfies $n=2h-1$, so $h\ge1$.
The hypothesis therefore makes $G$ bridgeless, and $N$ is connected.
Deleting $e$ decreases rank by one.
Suppressing $u,v$ preserves rank and connectedness and creates no loop.
This proves the first assertion.
The suppression occurs wholly within each original shore.
It preserves both shore ranks, while the other $k-1$ cut edges remain unchanged.
This proves the second assertion.

To prove minimality, take any cyclic cut of size $j$ in $G'$.
Lift its shore partition when restoring each suppressed vertex.
If the endpoints of its smoothed edge lie on the same shore, put the restored vertex on that shore.
If they lie on opposite shores, put it on either shore.
The corresponding two-edge path then has the same number of cut edges as the smoothed edge.
Every cycle contained in either shore lifts to a cycle on that shore.
Finally restore $e$, which increases the boundary size by at most one.
This gives a cyclic cut of size at most $j+1$ in $G$.
Hence $j\ge k-1$, proving the third assertion.

Suppression gives a bijection between the cycles of $N$ and those of $G'$.
The cycles of $G$ containing $e$ correspond bijectively to the simple $u$--$v$ paths in $N$.
This proves the final identity.
\end{proof}

\begin{lemma}[Uniform path bound after one edge deletion]
\label{lem:path-cut-deletion-path-bound}
Let $G$ range over connected loopless cubic multigraphs of rank $r$ with no cyclic cut of size at most four.
For every edge $e=uv$, uniformly in $G$ and $e$,
\[
 P_{uv}(G-e)\le(1/2+o(1))2^{r-1}
 \qquad(r\to\infty).
\]
\end{lemma}

\begin{proof}
The graph $G$ is bridgeless by the bridge observation in Lemma~\ref{lem:path-minimum-cyclic-cut-smoothing}.
Thus $G-e$ is a connected normalized two-pole of rank $r-1$, with terminals $u,v$.
Its cubic completion obtained by restoring $e$ is $G$.
For each fixed $\varepsilon>0$, Theorem~\ref{thm:high-density-pair-small-cyclic-cut} excludes path density at least $1/2+\varepsilon$ once $r-1\ge B_{2,\varepsilon}$.
The threshold is independent of $G$ and $e$, which proves the uniform estimate.
\end{proof}

\begin{theorem}[A relative deficit at every fixed minimum cut order]
\label{thm:path-fixed-cyclic-order-deficit}
For every fixed integer $k\ge5$, connected loopless cubic multigraphs of rank $r$ with $\lambda_c(G)=k$ satisfy
\[
 c(G)\le\left(\frac12-2^{-k-1}+o_k(1)\right)2^r
 \qquad(r\to\infty),
\]
uniformly over these graphs.
\end{theorem}

\begin{proof}
The five-cut proof gives the required base case with its explicit relative gap.
Indeed, for $\lambda_c(G)=5$, write the shore ranks as $h\le l$.
Proposition~\ref{prop:five-cut-universal-deficit} and the internal-cycle estimate in Theorem~\ref{thm:sharp-order-five-cut-error} give
\[
 \frac{c(G)}{2^{h+l}}\le\frac{31}{4}+o(1).
\]
Since $r=h+l+4$, this is
\[
 \frac{c(G)}{2^r}\le\frac{31}{64}+o(1)
 =\frac12-2^{-6}+o(1).
\]

Suppose the theorem holds for $k-1\ge5$.
Let $\lambda_c(G)=k$ and $\beta(G)=r$.
At sufficiently large rank, $G$ has girth at least five.
Otherwise a cycle of length at most four has a cyclic complement and gives an excluded cut of size at most four.
Apply Lemma~\ref{lem:path-minimum-cyclic-cut-smoothing} to one edge of a minimum cyclic cut.
The resulting graph $G'$ has rank $r-1$ and minimum cyclic cut order $k-1$.
The induction hypothesis gives
\[
 c(G')\le\left(\frac12-2^{-k}+o_k(1)\right)2^{r-1}.
\]
Lemma~\ref{lem:path-cut-deletion-path-bound} gives
\[
 P_{uv}(G-e)\le(1/2+o(1))2^{r-1}.
\]
Adding these estimates through the exact cycle identity yields
\[
 c(G)\le\left(1-2^{-k}+o_k(1)\right)2^{r-1}
 =\left(\frac12-2^{-k-1}+o_k(1)\right)2^r.
\]
All errors are uniform at each fixed induction step.
This completes the induction.
\end{proof}

\begin{corollary}[Sharp error for every fixed cyclic cut order]
\label{cor:path-all-fixed-cyclic-cuts-sharp}
For every fixed integer $K\ge1$, there is a constant $C_K$ such that every connected loopless cubic multigraph $G$ of rank $r$ with a cyclic cut of size at most $K$ satisfies
\[
 c(G)\le2^{r-1}+C_K r2^{r/2}.
\]
For all sufficiently large ranks, if $G$ has no cyclic cut of size at most four, then the stronger uniform estimate holds:
\[
 c(G)\le\left(\frac12-2^{-K-1}+o_K(1)\right)2^r<2^{r-1}.
\]
The second assertion is vacuous when $K\le4$.
\end{corollary}

\begin{proof}
For a cyclic cut of size at most four, use the established sharp four-cut theorem.
Otherwise the minimum cyclic cut order is one of $5,\ldots,K$.
Apply Theorem~\ref{thm:path-fixed-cyclic-order-deficit} to these finitely many orders.
The ranks below the resulting thresholds are absorbed into $C_K$, using $c(G)\le2^r-1$.
\end{proof}

Consequently, any sequence of above-half cubic multigraphs with ranks tending to infinity either has small cyclic cuts or has minimum cyclic cut order tending to infinity, in the extended sense.
After the sharp four-cut class is removed, a counterexample sequence to the unrestricted sharp-order bound must have cyclic connectivity tending to infinity.
The unrestricted bound $O(r2^{r/2})$ remains open.

For sufficiently large cubic rank, the minimum cyclic cut order is finite and logarithmic in the rank.
Indeed, the breadth-first tree bound gives a shortest cycle of length
\[
 g\le2\log_2(2r-2)+2<r.
\]
A shortest cycle is induced at these ranks.
Deleting its vertices leaves rank at least $r-g>0$, by the cubic degree count.
It therefore defines a cyclic cut of size $g$.
Thus the remaining unbounded-order regime still satisfies
\[
 \lambda_c(G)\le2\log_2(2r-2)+2.
\]

\subsection{A finite recurrence for a growing cut order}

For an integer $s$, define the edge-deletion path excess
\[
 q(s)=\sup_{G,e}\left(2^{1-s}P_{uv}(G-e)-\frac12\right)_+,
\]
where $G$ has rank $s$, is connected loopless cubic, has no cyclic cut of size at most four, and $e=uv$.
Here and below, cubic graphs may have parallel edges.
Set a supremum over an empty class equal to zero.
Lemma~\ref{lem:path-cut-deletion-path-bound} gives $q(s)\to0$.
Define the nonnegative five-cut base error by
\[
 a_5(s)=\sup_{\substack{\beta(G)=s\\\lambda_c(G)=5}}
 \left(2^{-s}c(G)-\frac{31}{64}\right)_+.
\]
The uniform five-cut estimate gives $a_5(s)\to0$.

\begin{proposition}[Finite cut-order recurrence]
\label{prop:path-finite-cut-recurrence}
Let $G$ be connected loopless cubic, with rank $r$ and minimum cyclic cut order $k\ge5$.
Assume $r-k+5\ge10$.
Then
\begin{align*}
 \frac{c(G)}{2^r}
 \le{}&\frac12-2^{-k-1}
 +2^{-(k-5)}a_5(r-k+5)\\
 &+\sum_{i=0}^{k-6}2^{-i-1}q(r-i).
\end{align*}
The sum is empty when $k=5$.
\end{proposition}

\begin{proof}
Apply Lemma~\ref{lem:path-minimum-cyclic-cut-smoothing} repeatedly, each time to an edge of a minimum cyclic cut.
After $j=k-5$ steps, the rank is $r-j$ and the minimum cyclic cut order is five.
Every intermediate rank is at least ten, so the absence of a cyclic cut of size at most four guarantees the needed girth condition.
The exact cycle identity and the definition of $q$ give at step $i$
\[
 c(G_i)\le c(G_{i+1})+(1/2+q(r-i))2^{r-i-1}.
\]
Sum these identities for $0\le i<j$ and use
\[
 c(G_j)\le\left(\frac{31}{64}+a_5(r-j)\right)2^{r-j}.
\]
The geometric sum of the half terms gives the stated formula.
\end{proof}

For a sequence with $k=k(r)\ge5$, the same argument gives a strict half-bound deficit whenever
\[
 r-k\longrightarrow\infty,
 \qquad
 2^k\max_{r-k+6\le s\le r}q(s)\longrightarrow0.
\]
For $k=5$, the maximum condition is unnecessary.
The base-error term is then negligible relative to $2^{-k-1}$, as is the weighted path-excess sum.

There is also an explicit sufficient threshold in terms of additive path stability.
Choose $R_5$ so that $a_5(s)\le1/256$ for $s\ge R_5$.
Set $\varepsilon=2^{-k-3}$ and
\[
 B_{2,\varepsilon}=3\bigl(9\lceil C_\varepsilon\rceil+25\bigr).
\]
If
\[
 r-k+5\ge\max\{10,R_5,B_{2,2^{-k-3}}\},
\]
then the path theorem gives $q(s)\le2^{-k-3}$ at every rank in the recurrence.
The two error terms are each at most one quarter of $2^{-k-1}$.
Consequently
\[
 c(G)\le\left(\frac12-2^{-k-2}\right)2^r.
\]
This condition permits a cut order that grows sufficiently slowly with rank.
Its quantitative limitation is the dependence of $C_\varepsilon$ on $\varepsilon$.
